\documentclass[english]{../../../unipd-notes}
\unipdsetup{
  course = {Computer Engineering Fundamentals},
  short-course = {Engineering Fundamentals},
  document-type = {Lecture notes},
  author = {Simone Siega},
  date = {3 August 2026}
}
\begin{document}
\makecoursefrontmatter
\begin{coursepreface}
These notes follow the course syllabus and complement the lectures with worked examples and self-assessment exercises.
\end{coursepreface}
\begin{revisionhistory}
  \revision{0.1.0}{10 July 2026}{Created the initial course outline}
  \revision{0.2.0}{12 July 2026}{Added the processor organization chapter}
  \revision{0.3.0}{15 July 2026}{Expanded the instruction-cycle examples}
  \revision{0.4.0}{18 July 2026}{Added the first memory hierarchy diagrams}
  \revision{0.5.0}{20 July 2026}{Revised the cache mapping explanation}
  \revision{0.6.0}{22 July 2026}{Added exercises on address translation}
  \revision{0.7.0}{24 July 2026}{Expanded the input and output chapter}
  \pagebreak
  \revision{0.8.0}{26 July 2026}{Added the interrupt-handling examples}
  \revision{0.9.0}{28 July 2026}{Reviewed terminology and cross-references}
  \revision{1.0.0}{1 August 2026}{Published the introductory lectures}
  \revision{1.1.0}{2 August 2026}{Added the memory hierarchy summary}
  \revision{1.2.0}{3 August 2026}{Corrected examples and improved navigation}
\end{revisionhistory}
\makecoursemainmatter
\chapter{Introduction to computer systems}
\begin{figure}[htbp]
  \centering\fbox{\rule{0pt}{20mm}\rule{35mm}{0pt}}
  \caption{Functional organization of a computer}
\end{figure}
\printcourselists
\courseappendices
\chapter{Exam reference sheet}
This appendix summarizes binary prefixes, base conversions, and characteristic memory hierarchy timings.
\end{document}
